\par\medskip\noindent\textbf{Cómo leer la salida.}
{\small\setlength{\tabcolsep}{4pt}\renewcommand{\arraystretch}{1.15}
\begin{xltabular}{\linewidth}{@{}>{\raggedright\arraybackslash}p{0.24\linewidth}>{\raggedleft\arraybackslash}p{0.15\linewidth}>{\raggedright\arraybackslash}X@{}}
\toprule Indicador & Valor & Interpretación \\ \midrule \endfirsthead
\toprule Indicador & Valor & Interpretación \\ \midrule \endhead
\bottomrule \endlastfoot
$p$ por BIC, $k$ & 3; $k=2$ & Rezagos del VAR en niveles y, por tanto, de las diferencias en el VECM. \\
Autovalores & 0,0631; 0,0050 & Correlaciones canónicas al cuadrado; uno grande y otro casi cero sugieren un vector de cointegración. \\
Traza $r\le0$ & 33,05 (VC 15,49) & Se rechaza que no haya cointegración. \\
Traza $r\le1$ & 2,35 (VC 3,84) & No se rechaza que haya como máximo un vector: $r=1$. \\
Máximo autovalor & 30,70 / 2,35 & Prueba alternativa con la misma conclusión. \\
$\bm\beta$ normalizado & $(1;\ -1{,}112)$ & Relación de largo plazo $B_t-1{,}112\,W_t$ estacionaria. \\
$\alpha_{\text{Brent}}$ & 0,034 ($t=0{,}40$) & El Brent no reacciona al desequilibrio: es el precio ``líder''. \\
$\alpha_{\text{WTI}}$ & 0,166 ($t=2{,}03$) & El WTI corrige el desequilibrio: se ajusta hacia el Brent cada mes. \\
Persistencia $1+\beta'\alpha$ & 0,849 & Parte del desequilibrio que persiste al mes siguiente. \\
Vida media & 4,2 meses & Coherente con la de Engle-Granger. \\
LR $\beta=(1,-1)$ & 22,39; $p\approx2\times10^{-6}$ & Se rechaza que los precios se muevan uno a uno en el largo plazo en toda la muestra. \\
\end{xltabular}}
